% theory.tex

\chapter{相关理论与技术}
\echapter{Related Theory and Techniques}

\section{时序数据预测任务定义}
\esection{Definition of Time Series Forecasting}

时间序列预测任务的目标是根据历史观测数据估计未来一段时间内的取值。设多变量时间序列为 $\mathbf{X}=\{\mathbf{x}_1,\mathbf{x}_2,\cdots,\mathbf{x}_T\}$，其中 $\mathbf{x}_t \in \mathbb{R}^{C}$ 表示第 $t$ 个时间步的 $C$ 个变量观测值。常见的多步预测任务通常给定长度为 $L$ 的历史输入序列：
\begin{equation}
    \mathbf{X}_{in} = [\mathbf{x}_{t-L+1}, \mathbf{x}_{t-L+2}, \cdots, \mathbf{x}_{t}],
\end{equation}
要求模型预测未来长度为 $H$ 的序列：
\begin{equation}
    \hat{\mathbf{Y}} = [\hat{\mathbf{x}}_{t+1}, \hat{\mathbf{x}}_{t+2}, \cdots, \hat{\mathbf{x}}_{t+H}].
\end{equation}
在本文实验中，输入长度统一设置为 $L=96$，预测长度 $H$ 分别取 96、192、336 和 720。
这种设置对应长序列预测中较常见的多步预测场景。与只预测下一时刻的单步预测相比，多步预测需要模型同时处理短期局部变化和较长跨度上的趋势延续。当预测步长较短时，模型更容易利用最近观测值中的局部形态；当预测步长增加时，未来序列与输入窗口之间的直接对应关系减弱，误差也更容易在时间维度上放大。因此，长预测任务不仅要求模型拟合局部波动，还要求模型保留对周期结构和整体趋势的描述能力。本文将输入长度固定为 96，是为了在不同模型变体之间保持一致的历史信息范围；同时设置多个预测长度，则可以观察模型在短、中、长不同预测跨度下的表现差异，而不是只根据单一预测步长判断模型优劣。

从数据形式看，时序数据预测既可以是单变量预测，也可以是多变量预测。本文采用的公开数据集大多为多变量序列，例如 Weather 数据集包含多个气象指标，Traffic 数据集包含多个道路传感器的占用率。对于多变量预测，模型不仅需要利用每个变量自身的历史变化，还可能需要利用变量之间的相关性。
时序数据预测的困难主要来自三个方面。首先，预测步长增加时，模型需要更好地把握趋势和周期结构，否则误差容易逐渐放大。其次，不同数据集的时间模式差异较大，同一模型在不同场景下可能表现不一致。最后，多变量序列中的变量数可能较多，如果直接进行全变量交互，模型计算成本会明显增加。因此，本文选择 TimeMixer 这类轻量级多尺度模型作为基础，并重点分析其结构改进的效果。

\section{时间序列分解方法}
\esection{Time Series Decomposition}

时间序列通常可以被看作由多个成分共同组成，其中较常见的划分是趋势项、季节项和随机噪声项。趋势项描述序列在较长时间范围内的整体变化方向，季节项描述具有周期性的波动，噪声项则表示难以建模的随机扰动。虽然真实数据很难完全满足简单分解假设，但这种思想可以帮助模型从不同角度理解序列结构。
在预测任务中，不同成分承担的作用并不相同。趋势项通常影响较长预测步长下的整体水平，如果趋势估计出现偏差，模型输出可能在整个预测区间内系统性偏高或偏低；季节项更多对应日周期、周周期或设备运行周期等重复模式，能够帮助模型在未来时间段中恢复相似波动；噪声项虽然难以准确预测，但如果模型把噪声误认为趋势或周期，反而可能降低泛化能力。因此，分解方法的价值不只在于把序列拆开，还在于为模型提供一种更清晰的建模顺序：先估计较平滑的结构，再处理剩余的局部变化。

在深度学习预测模型中，分解模块常被用作预处理或中间结构。假设输入序列为 $\mathbf{X}$，分解过程可表示为：
\begin{equation}
    \mathbf{X} = \mathbf{S} + \mathbf{T},
\end{equation}
其中 $\mathbf{S}$ 表示季节项或残差项，$\mathbf{T}$ 表示趋势项。滑动平均是一种常见的趋势估计方法。给定窗口大小 $k$，趋势项可以由邻域内时间点的平均值得到。原始 TimeMixer 采用的移动平均分解方式结构简单、计算稳定，但固定窗口也意味着模型对不同周期长度的适应能力有限。
滑动平均的优点在于形式直观、参数量低，并且不会引入额外复杂结构。对于包含一定噪声但整体变化较平缓的序列，滑动平均能够起到平滑作用，使模型更容易关注主要趋势。然而，窗口大小一旦确定，分解结果也会随之固定。较小窗口能够保留短期变化，但对长期趋势的平滑不足；较大窗口能够得到更平滑的趋势估计，却可能削弱局部突变和短周期波动。实际数据往往同时包含多种时间尺度，例如交通数据中既有早晚高峰，也有局部拥堵；太阳能数据中既有昼夜周期，也受天气扰动影响。单一固定窗口很难同时覆盖这些变化。

本文提出的多核门控分解模块仍然保留移动平均的基本思想，但不再只使用单一窗口。它并行计算多个窗口下的趋势估计，并根据每个样本的全局均值特征生成融合权重。这种设计可以看作是在传统分解方法和轻量级神经网络之间做折中：一方面保留分解结构的直观性，另一方面通过样本级门控权重增强一定的自适应能力。

\section{多尺度建模思想}
\esection{Multi-scale Modeling}

多尺度建模是处理复杂时间序列的一种常见思路。实际序列往往同时包含不同时间尺度的信息：短时间尺度上可能出现突变和局部波动，较长时间尺度上则体现整体趋势和周期规律。单一时间分辨率可能难以同时兼顾这些信息，因此许多模型会构建不同尺度的表示。

一种直接做法是对原始序列进行下采样，得到多个时间分辨率的序列。例如，原始序列保留完整时间步，下一级序列每隔若干点采样或池化一次，再下一级序列进一步降低时间分辨率。这样形成的序列金字塔能够为模型提供不同尺度的信息。高分辨率序列更适合描述局部细节，低分辨率序列更适合描述长期变化。
从信息分布角度看，高分辨率尺度保留了更多相邻时间点之间的细节，适合捕捉突发变化、短期震荡和局部峰值；低分辨率尺度通过聚合或降采样压缩了时间维度，更容易呈现较长时间范围内的趋势和周期。二者并不是简单替代关系，而是从不同时间粒度观察同一序列。对于长序列预测，如果只依赖高分辨率输入，模型需要在较长时间维度上处理大量细节；如果只依赖低分辨率输入，又可能丢失对局部变化的响应能力。多尺度结构的意义在于让模型在不同时间粒度之间分工，并在后续混合过程中重新整合这些信息。

但是，下采样本身也会带来信息损失。如果下采样方式过于简单，重要的局部信息可能在进入后续网络前已经被削弱。平均池化虽然平滑稳定，但无法根据任务自动选择关键时间点。本文在 TimeMixer 的多尺度构建阶段引入深度可分离卷积下采样，希望让下采样过程具备一定学习能力，在保持轻量化的同时减少固定池化带来的局限。
多尺度建模也适合与轻量化预测模型结合。一方面，下采样后的低分辨率序列长度更短，能够降低后续跨尺度混合的计算压力；另一方面，不同尺度表示为模型提供了结构化归纳偏置，使模型不必完全依赖大规模参数去学习所有时间模式。TimeMixer 正是利用这一点，在不使用复杂注意力机制的情况下完成多尺度分解和混合。本文的改进也延续这一思路：增强多尺度构建、分解和变量处理环节，但不改变模型以轻量化结构处理长序列预测的基本方向。

\section{原始TimeMixer模型结构}
\esection{Original TimeMixer Architecture}

TimeMixer\cite{wang2023timemixer} 是一种基于多尺度分解混合思想的时间序列预测模型。它没有采用复杂的自注意力结构，而是通过多尺度序列构建、过去可分解混合模块和未来多尺度预测融合模块完成预测。原始模型整体结构如图 \ref{fig:timemixer_overall} 所示。

\begin{figure}[htbp]
    \centering
    \includegraphics[width=0.82\textwidth]{thesis/fig_timemixer_original_framework.png}
    \caption{TimeMixer原始模型结构示意图（改绘自文献\cite{wang2023timemixer}）}
    \label{fig:timemixer_overall}
\end{figure}

从整体流程看，TimeMixer 的输入为长度为 $P$、变量数为 $C$ 的历史序列 $\mathbf{X}\in\mathbb{R}^{P\times C}$，输出为长度为 $F$ 的未来预测序列 $\hat{\mathbf{Y}}\in\mathbb{R}^{F\times C}$。模型首先通过下采样得到一组不同时间分辨率的历史序列，再对这些序列进行嵌入和多层 PDM 处理，最后由 FMM 汇总不同尺度上的预测结果。该流程可以概括为表 \ref{tab:timemixer_workflow}。

\begin{table}[htbp]
    \centering
    \caption{原始TimeMixer模型工作流程}
    \label{tab:timemixer_workflow}
    \resizebox{\textwidth}{!}{%
    \begin{tabular}{@{}C{0.16\textwidth}L{0.25\textwidth}L{0.31\textwidth}L{0.24\textwidth}@{}}
        \toprule
        \multicolumn{1}{c}{阶段} & \multicolumn{1}{c}{输入与输出} & \multicolumn{1}{c}{主要操作} & \multicolumn{1}{c}{作用} \\
        \midrule
        多尺度序列构建 & 输入原始历史序列，输出多个时间分辨率的序列集合 & 通过平均池化逐级下采样，并对各尺度序列进行嵌入 & 将复杂时间变化拆分到细粒度和粗粒度尺度中 \\
        过去可分解混合 & 输入多尺度历史表示，输出混合后的多尺度表示 & 对每个尺度分解季节项和趋势项，分别进行跨尺度混合 & 从历史观测中提取细节波动和宏观趋势信息 \\
        未来多尺度预测融合 & 输入最终多尺度表示，输出未来预测序列 & 每个尺度设置独立预测器，并将多个尺度预测结果相加融合 & 利用不同尺度预测器的互补能力形成最终预测 \\
        \bottomrule
    \end{tabular}}
\end{table}

\subsection{多尺度序列构建}
\esubsection{Multi-scale Sequence Construction}

TimeMixer 首先对输入序列进行多尺度处理。原文指出，不同采样尺度下的时间序列往往呈现不同性质：细尺度更容易保留局部细节和短期波动，粗尺度则更突出长期变化和宏观趋势。因此，TimeMixer 没有只在单一时间分辨率上建模，而是先通过平均池化生成多尺度历史观测。给定原始输入 $\mathbf{X}$，模型通过下采样得到多个不同长度的序列：
\begin{equation}
    \{\mathbf{X}^{0}, \mathbf{X}^{1}, \cdots, \mathbf{X}^{M}\},
\end{equation}
其中 $\mathbf{X}^{0}$ 为原始尺度，$\mathbf{X}^{i}$ 表示第 $i$ 个下采样尺度，$M$ 为下采样层数。随着尺度加深，序列长度逐渐缩短，时间分辨率降低。

更具体地说，若原始输入长度为 $P$，第 $m$ 个尺度的序列长度约为 $\lfloor P/2^m\rfloor$，可表示为：
\begin{equation}
    \mathbf{X}_{ms}=\{\mathbf{x}_0,\mathbf{x}_1,\cdots,\mathbf{x}_M\},\quad
    \mathbf{x}_m\in \mathbb{R}^{\lfloor P/2^m\rfloor\times C}.
\end{equation}
其中 $\mathbf{x}_0$ 对应原始输入，保留最完整的时间细节；$\mathbf{x}_M$ 对应最粗尺度，更强调长时间范围内的变化轮廓。得到多尺度序列后，模型通过嵌入层将其投影到深层特征空间，形成后续 PDM 模块的输入表示。

这一设计的关键意义在于：多尺度序列不是简单地复制输入，而是为模型提供了不同时间粒度下的观察角度。对于局部波动明显的数据，细尺度序列可以保留较多短期变化；对于长期趋势明显的数据，粗尺度序列能够减少局部噪声干扰，使趋势更清晰。原始 TimeMixer 默认下采样方式为平均池化。平均池化的优点是实现简单且不会引入额外参数，但它对局部窗口内的所有时间点进行等权平均。对于时序数据预测任务，这种方式能够平滑噪声，但也可能削弱局部突变。本文的第一个改进模块正是围绕这一环节展开。

\subsection{过去可分解混合模块}
\esubsection{Past-Decomposable-Mixing Block}

过去可分解混合模块是 TimeMixer 的核心部分。原文认为，即使经过下采样，粗尺度序列中仍然可能同时包含趋势变化和季节波动；如果直接把多尺度序列作为整体混合，模型难以区分不同成分的性质。因此，PDM 先对每个尺度的表示进行序列分解，再分别混合季节项和趋势项。对于第 $l$ 层 PDM 中第 $m$ 个尺度的输入 $\mathbf{x}_m^{l}$，分解过程可表示为：
\begin{equation}
    \mathbf{s}_m^{l},\mathbf{t}_m^{l}=\mathrm{SeriesDecomp}(\mathbf{x}_m^{l}),
\end{equation}
其中 $\mathbf{s}_m^{l}$ 为季节项或局部波动成分，$\mathbf{t}_m^{l}$ 为趋势项。分解操作通常基于移动平均，将较平滑的部分作为趋势，将剩余部分作为季节或残差。

TimeMixer 对季节项采用自底向上的混合方式。高分辨率序列保留更多局部变化信息，原文将其理解为更细粒度的季节性细节。因此，季节项从高分辨率向低分辨率传递，有助于将细粒度波动逐步融入较粗尺度表示。对于相邻两个尺度，季节项混合可以写成：
\begin{equation}
    \mathbf{s}_m^{l}=\mathbf{s}_m^{l}+\mathrm{BottomUpMixing}(\mathbf{s}_{m-1}^{l}),\quad m=1,\cdots,M.
\end{equation}
其中 $\mathrm{BottomUpMixing}(\cdot)$ 可由沿时间维度的线性层和 GELU 激活构成，其输入长度对应较细尺度，输出长度对应较粗尺度。通过这种方式，低层细尺度中的局部周期和短期波动被逐级传递到高层尺度中。

与季节项相反，趋势项主要反映较长时间范围内的变化方向。细尺度中包含较多局部波动，这些波动对宏观趋势估计可能形成干扰；粗尺度经过下采样后更容易呈现整体趋势。因此，TimeMixer 对趋势项采用自顶向下的混合方式，让低分辨率中的宏观趋势信息回传到高分辨率表示中：
\begin{equation}
    \mathbf{t}_m^{l}=\mathbf{t}_m^{l}+\mathrm{TopDownMixing}(\mathbf{t}_{m+1}^{l}),\quad m=M-1,\cdots,0.
\end{equation}
其中 $\mathrm{TopDownMixing}(\cdot)$ 同样由时间维度上的线性映射和非线性激活实现，但其信息方向与季节项相反，由粗尺度指向细尺度。经过季节项和趋势项的分别混合后，PDM 将两类成分重新组合，并通过前馈网络进行通道方向的信息交互和残差更新。简化地看，第 $l$ 层 PDM 的输出可表示为：
\begin{equation}
    \mathbf{X}^{l+1}=\mathbf{X}^{l}+\mathrm{FeedForward}\big(\mathrm{S\mbox{-}Mix}(\mathbf{S}^{l})+\mathrm{T\mbox{-}Mix}(\mathbf{T}^{l})\big).
\end{equation}
其中 $\mathrm{S\mbox{-}Mix}(\cdot)$ 表示季节项自底向上混合，$\mathrm{T\mbox{-}Mix}(\cdot)$ 表示趋势项自顶向下混合。这样的双向分解混合结构是 TimeMixer 区别于普通多尺度拼接或直接混合方法的关键。

\subsection{未来多尺度预测融合模块}
\esubsection{Future-Multipredictor-Mixing Block}

经过多层 PDM 后，模型获得多个尺度的历史表示：
\begin{equation}
    \mathbf{X}^{L}=\{\mathbf{x}_0^{L},\mathbf{x}_1^{L},\cdots,\mathbf{x}_M^{L}\}.
\end{equation}
由于不同尺度中保留的信息不同，它们在未来预测阶段也具有不同能力。高分辨率尺度更关注短期细节，粗分辨率尺度更关注宏观趋势。TimeMixer 因此没有只使用最后一个尺度或简单拼接所有尺度，而是在 FMM 中为每个尺度设置独立预测器，再将多个预测器的输出相加得到最终结果：
\begin{equation}
    \hat{\mathbf{Y}}_m=\mathrm{Predictor}_m(\mathbf{x}_m^{L}),\quad
    \hat{\mathbf{Y}}=\sum_{m=0}^{M}\hat{\mathbf{Y}}_m.
\end{equation}
其中 $\hat{\mathbf{Y}}_m\in\mathbb{R}^{F\times C}$ 表示第 $m$ 个尺度给出的未来预测，$\hat{\mathbf{Y}}$ 为最终预测结果。每个预测器通常先沿时间维度将该尺度的历史长度映射到预测长度 $F$，再将深层特征投影回原始变量维度 $C$。

FMM 的思想可以理解为多尺度预测器集成。单个尺度预测器只能利用某一时间分辨率下的历史表示，可能偏向局部波动或长期趋势中的某一方面；多个尺度预测器相加后，可以同时利用细尺度和粗尺度的互补信息。原始 TimeMixer 的这种设计兼顾了表达能力和效率。它通过多尺度混合实现了类似“分而治之”的建模方式，同时避免了自注意力在较长输入窗口下可能带来的较高计算开销。由于原始结构本身已经较强，本文实验中 M0 基线仍然获得了较多最优次数，这一点在第 4 章实验分析中会进一步说明。

\section{原始模型不足分析}
\esection{Limitations of the Original Model}

综合上述结构可以看出，TimeMixer 的优势在于多尺度结构清晰、计算效率较高、分解混合思想适合时间序列预测。但从改进角度看，它仍有三个环节值得进一步优化。

首先，多尺度构建中的平均池化下采样缺乏学习能力。对于较平滑的数据，平均池化可能有利于抑制噪声；但对于包含突发变化的数据，它可能平滑掉有用信息。其次，固定窗口分解难以同时适应不同周期。不同数据集和不同预测长度下，最合适的趋势窗口可能并不相同。最后，通道独立建模虽然提高了效率，却弱化了变量之间的信息交互。对于 Weather、Traffic 等多变量数据集，这一点可能影响模型充分利用变量间关系。

因此，本文第 3 章分别从下采样、分解和变量交互三个角度提出改进模块。需要强调的是，这三个模块并不一定简单叠加后就能得到最优结果。本文通过 8 个变体的完整消融实验，对各模块的独立贡献和组合效果进行了分析。

\section{本章小结}
\esection{Summary}

本章首先给出了多变量多步时序预测任务的基本定义，说明了输入长度、预测长度和变量维度在预测问题中的含义；随后介绍了时间序列分解方法和多尺度建模思想，分析了趋势项、季节项以及不同时间分辨率表示在长序列预测中的作用。在此基础上，本章进一步梳理了原始 TimeMixer 模型的多尺度序列构建、过去可分解混合和未来多尺度预测融合结构，并从平均池化下采样、固定窗口分解和通道独立处理三个方面总结了原模型可改进之处。这些内容为第 3 章提出深度可分离卷积下采样、多核门控分解和通道混合模块提供了理论和结构基础。
